\chapter*{Przedmowa}
\thispagestyle{empty}

Osobom postronnym często się wydaje, że badania matematyczne nie mają związku z otaczającą nas rzeczywistością. Jest to konsekwencja wysoce sformalizowanego języka, którym współczesnie posługuje się matematyka. Za tym formalizmem kryje się jednak bardzo bogaty świat matematycznych intuicji, bezpośrednio inspirowany nie tylko konkretnymi problemami nauki i techniki, ale przede wszystkim doświadczeniem i wyobrażeniami z codziennego życia. Matematycy rozwijają swoją intuicję wymyślając i analizując liczne przykłady. W przykładach tych poszukują powtarzających się prawidłowości. Jeśli taką prawidłowość uda się wykryć i nie udaje się znaleźć przykładu, w którym prawidłowość nie zachodzi, formułuje się hipotezę, że prawidłowość taka zachodzi zawsze. Hipoteza nie musi być i często nie jest w pełni prawidłowym opisem rzeczywistości, choć przynajmniej ziarno prawdy często w niej tkwi.

W odróżnieniu od wielu innych nauk, matematycy odrzucają intuicję i eksperyment jako metodę rozstrzygającą co jest prawdą. Hipotezę trzeba uzasadnić poprzez wyprowadzenie jej drogą logicznego rozumowania z uznawanych za oczywiste faktów. Rozumowanie takie nazywa się dowodem, a hipotezę, dla której udało się podać dowód, twierdzeniem. W tym miejscu pojawia się sformalizowany język. Przeprowadzenie matematycznego dowodu wymaga bowiem stosownej teorii. W teorii takiej ustala się najpierw podstawowe pojęcia, czyli tak zwane pojęcia pierwotne teorii oraz fakty uznawane za oczywiste, nazywane aksjomatami teorii. Wszystkie inne pojęcia teorii definiowane są w oparciu o pojęcia pierwotne, a twierdzenia dowodzone są w oparciu o przyjęte aksjomaty. Wysoki stopień sformalizowania języka matematycznego dowodu ma uniemożliwić, a przynajmniej utrudnić, świadome lub nieświadome przemycenie do dowodu fałszywego argumentu.

Tytuły przepisów po których następuje gwiazdka oznaczają dania mięsne.

\bigskip
\begin{flushright}
    \textit{"In these days the angel of topology and the devil of abstract algebra fight for the soul of every individual discipline of mathematics."} \\
    \textit{- Hermann Weyl}
\end{flushright}

\newpage